\section{Digital Signal Processing (DSP) - Session
3}\label{digital-signal-processing-dsp---session-3}

\begin{center}\rule{0.5\linewidth}{0.5pt}\end{center}

\subsection{Discrete-Time Signals and
Systems}\label{discrete-time-signals-and-systems}

\subsubsection{Analysis of Discrete-Time Linear Time-Invariant
Systems}\label{analysis-of-discrete-time-linear-time-invariant-systems}

We aim to design \textbf{Linear Time-Invariant (LTI) systems} that are
both \textbf{stable} and \textbf{causal}.

There are two basic methods for analyzing the response of a linear
system to a given input signal:

\begin{itemize}
\tightlist
\item
  \textbf{Direct Solution:} Based on the direct solution of the
  input-output (difference) equation for the system.
\item
  \textbf{Decomposition:} Decomposing the input signal into a sum of
  elementary signals, selected so that the response of the system to
  each component is easily determined.
\end{itemize}

\subsubsection{Resolution of a Discrete-Time Signal into
Impulses}\label{resolution-of-a-discrete-time-signal-into-impulses}

A discrete-time signal \(x(n)\) can be resolved into a weighted sum of
unit sample (impulse) sequences. Each multiplication of the signal
\(x(n)\) by a unit impulse at a specific delay \(k\) picks out the
single value \(x(k)\).
\[x(n) = \sum_{k=-\infty}^{\infty} x(k)\delta(n-k)\]

\subsubsection{Response of LTI Systems: The Convolution
Sum}\label{response-of-lti-systems-the-convolution-sum}

Given a LTI system \(\mathcal{T}\):

\begin{itemize}
\tightlist
\item
  The system's response to a unit sample \(\delta(n)\) is the
  \textbf{impulse response}. \[h(n) = \mathcal{T}[\delta(n)]\]
\item
  By the property of time-invariance, the response to a delayed impulse
  \(\delta(n-k)\) is \(h(n-k)\).
\item
  By the property of linearity (superposition), the total output
  \(y(n)\) to an arbitrary input \(x(n)\) is the sum of the weighted
  responses to each impulse component: \[
    y(n) = \mathcal{T}[x(n)] = \mathcal{T}\left[\sum_{k=-\infty}^{\infty} x(k)\delta(n-k)\right] = \sum_{k=-\infty}^{\infty} x(k)\mathcal{T}[\delta(n-k)] = \sum_{k=-\infty}^{\infty} x(k)h(n-k)
    \]
\item
  The impulse response \(h(n)\) completely characterizes the LTI system,
  as the output can be determined for any input using this
  transformation.
\end{itemize}

The operation \[
\sum_{k=-\infty}^{\infty} x(k)h(n-k)
\] is called the \textbf{convolution} of \(x(n)\) and \(h(n)\), often
denoted as

\[x(n) * h(n)\]

\subsubsection{Properties of
Convolution}\label{properties-of-convolution}

Convolution satisfies several mathematical laws:

\begin{itemize}
\tightlist
\item
  \textbf{Commutative Law:} \[x(n) * h(n) = h(n) * x(n)\]
\item
  \textbf{Associative Law:}
  \[[x(n) * h_1(n)] * h_2(n) = x(n) * [h_1(n) * h_2(n)]\]
\item
  \textbf{Distributive Law:}
  \[x(n) * [h_1(n) + h_2(n)] = x(n) * h_1(n) + x(n) * h_2(n)\]
\item
  \textbf{Identity Property:} \[x(n) * \delta(n) = x(n)\]
\item
  \textbf{Shifting Property:} \[x(n) * \delta(n-k) = x(n-k)\]
\end{itemize}

\subsubsection{Steps for Computing
Convolution}\label{steps-for-computing-convolution}

The process of evaluating the convolution sum involves four distinct
operations:

\begin{enumerate}
\def\labelenumi{\arabic{enumi}.}
\tightlist
\item
  \textbf{Folding:} Fold \(h(k)\) about \(k=0\) to obtain \(h(-k)\).
\item
  \textbf{Shifting:} Shift the folded sequence by \(n\) to the right (if
  \(n>0\)) or left (if \(n<0\)) to obtain \(h(n-k)\).
\item
  \textbf{Multiplication:} Multiply \(x(k)\) by \(h(n-k)\) to obtain the
  product sequence.
\item
  \textbf{Summation:} Sum all values of the product sequence to obtain
  the output value \(y(n)\).
\end{enumerate}

\begin{quote}
\textbf{\emph{Example 3.2}}

The impulse response of a linear time-invariant system is

\[
h(n) = \{1, \underset{\uparrow}{2}, 1, -1\}
\]

Determine the response of the system to the input signal

\[
x(n) = \{\underset{\uparrow}{1}, 2, 3, 1\}
\]

\emph{Solution}

\begin{enumerate}
\def\labelenumi{\arabic{enumi}.}
\tightlist
\item
  \textbf{Define Input Sequences}:
\end{enumerate}

\begin{itemize}
\tightlist
\item
  \textbf{Input signal:}

  \begin{itemize}
  \tightlist
  \item
    \(x(k) = \{\underset{\uparrow}{1}, 2, 3, 1\}, \text{ for } k = 0, 1, 2, 3\)
  \end{itemize}
\item
  \textbf{Impulse response:}

  \begin{itemize}
  \tightlist
  \item
    \(h(k) = \{1, \underset{\uparrow}{2}, 1, -1\}, \text{ for } k = -1, 0, 1, 2\)
  \end{itemize}
\end{itemize}

\begin{enumerate}
\def\labelenumi{\arabic{enumi}.}
\setcounter{enumi}{1}
\item
  \textbf{Convolution Sum Formula}:
  \[y(n) = \sum_{k=-\infty}^{\infty} x(k)h(n-k)\]
\item
  \textbf{Step-by-Step Calculation}:

  \begin{itemize}
  \tightlist
  \item
    \textbf{For \(n = -1\)}:

    \begin{itemize}
    \tightlist
    \item
      \(y(-1) = x(0)h(-1) = (1)(1) = 1\)
    \end{itemize}
  \item
    \textbf{For \(n = 0\)}:

    \begin{itemize}
    \tightlist
    \item
      \(y(0) = x(0)h(0) + x(1)h(-1) = (1)(2) + (2)(1) = 4\)
    \end{itemize}
  \item
    \textbf{For \(n = 1\)}:

    \begin{itemize}
    \tightlist
    \item
      \(y(1) = x(0)h(1) + x(1)h(0) + x(2)h(-1) = (1)(1) + (2)(2) + (3)(1) = 8\)
    \end{itemize}
  \item
    \textbf{For \(n = 2\)}:

    \begin{itemize}
    \tightlist
    \item
      \(y(2) = x(0)h(2) + x(1)h(1) + x(2)h(0) + x(3)h(-1) = (1)(-1) + (2)(1) + (3)(2) + (1)(1) = 8\)
    \end{itemize}
  \item
    \textbf{For \(n = 3\)}:

    \begin{itemize}
    \tightlist
    \item
      \(y(3) = x(1)h(2) + x(2)h(1) + x(3)h(0) = (2)(-1) + (3)(1) + (1)(2) = 3\)
    \end{itemize}
  \item
    \textbf{For \(n = 4\)}:

    \begin{itemize}
    \tightlist
    \item
      \(y(4) = x(2)h(2) + x(3)h(1) = (3)(-1) + (1)(1) = -2\)
    \end{itemize}
  \item
    \textbf{For \(n = 5\)}:

    \begin{itemize}
    \tightlist
    \item
      \(y(5) = x(3)h(2) = (1)(-1) = -1\)
    \end{itemize}
  \end{itemize}
\item
  \textbf{Final Response Sequence}:
  \[y(n) = \{1, \underset{\uparrow}{4}, 8, 8, 3, -2, -1\}\]
\end{enumerate}
\end{quote}

\begin{quote}
\textbf{\emph{Example 3.3}}

Determine the output \(y(n)\) of a relaxed linear time-invariant system
with impulse response

\[h(n) = a^n u(n), \quad |a| < 1\]

when the input is a unit step sequence, that is,

\[x(n) = u(n)\]

\emph{Solution}

\begin{enumerate}
\def\labelenumi{\arabic{enumi}.}
\tightlist
\item
  \textbf{Define the sequences:} Both the impulse response \(h(n)\) and
  the input signal \(x(n)\) are infinite-duration sequences.
\item
  \textbf{Select the convolution form:} The convolution is computed
  using the form where the input sequence \(x(k)\) is folded and
  shifted: \[y(n) = \sum_{k=-\infty}^{\infty} x(n-k)h(k)\]
\item
  \textbf{Evaluate for \(n < 0\):} When the shift \(n\) is negative,
  there is no overlap between the non-zero values of \(x(n-k)\) and
  \(h(k)\). Consequently, the product sequences consist of all zeros, so
  \textbf{\(y(n) = 0\) for \(n < 0\)}.
\item
  \textbf{Evaluate for \(n \ge 0\):} Compute the sum of the product
  sequences for successive values of \(n\):

  \begin{itemize}
  \tightlist
  \item
    For \(n=0\): \(y(0) = 1\).
  \item
    For \(n=1\): \(y(1) = 1 + a\).
  \item
    For \(n=2\): \(y(2) = 1 + a + a^2\).
  \end{itemize}
\item
  \textbf{Generalize the summation:} For any \(n > 0\), the output is
  the sum of a geometric series:
  \[y(n) = 1 + a + a^2 + \dots + a^n = \sum_{k=0}^{n} a^k\]
\item
  \textbf{Simplify the result:} Applying the geometric series formula,
  the output for \(n \ge 0\) is: \[y(n) = \frac{1 - a^{n+1}}{1 - a}\]
\item
  \textbf{Identify steady-state response:} Since \(|a| < 1\), as \(n\)
  approaches infinity, the output reaches a final value of:
  \[y(\infty) = \frac{1}{1 - a}\]
\end{enumerate}
\end{quote}

\subsubsection{Causality and Stability
Conditions}\label{causality-and-stability-conditions}

For an LTI system to be practically realizable and useful, it must
satisfy specific conditions regarding its impulse response:

\begin{itemize}
\tightlist
\item
  \textbf{Causality:} An LTI system is causal if and only if its impulse
  response is zero for negative time: \[h(n) = 0, \quad n < 0\]
\item
  \textbf{Stability (BIBO):} An LTI system is Bounded-Input
  Bounded-Output (BIBO) stable if and only if its impulse response is
  \textbf{absolutely summable}:
  \[S_h = \sum_{k=-\infty}^{\infty} |h(k)| < \infty\]
\end{itemize}

\subsubsection{Systems with Finite-Duration and Infinite-Duration
Impulse Response (FIR and
IIR)}\label{systems-with-finite-duration-and-infinite-duration-impulse-response-fir-and-iir}

\textbf{Finite-Duration Impulse Response (FIR)}

\[ h(n) = 0, \quad n < 0, n ≥ M \]

\[ y(n) = \sum_{k=0}^{M-1} h(k)x(n-k) \]

\begin{itemize}
\tightlist
\item
  \textbf{Description:} The impulse response is zero outside of a finite
  time interval.
\item
  \textbf{Interpretation:} The output \(y(n)\) at any time is a
  \textbf{weighted linear combination} of the most recent \(M\) input
  signal samples, specifically \(x(n), x(n-1), \dots, x(n-M+1)\).
\item
  \textbf{Memory:} The system acts as a \textbf{window} that only views
  the most recent \(M\) samples and ``forgets'' all prior inputs, giving
  it a \textbf{finite memory} of length \(M\).
\item
  \textbf{Implementation:} FIR systems are easily implemented directly
  as suggested by the convolution sum, requiring additions,
  multiplications, and a \textbf{finite number of memory locations}.
  They are typically realized using \textbf{nonrecursive} structures.
\end{itemize}

\textbf{Infinite-Duration Impulse Response (IIR)}

\[ y(n) = \sum_{k=0}^{\infty} h(k)x(n-k) \]

\begin{itemize}
\tightlist
\item
  \textbf{Description:} The impulse response has an infinite duration.
\item
  \textbf{Interpretation:} The output is a weighted linear combination
  of the present and \textbf{all past input samples}.
\item
  \textbf{Memory:} Because the output depends on the entire history of
  the input signal, the system is said to have \textbf{infinite memory}.
\item
  \textbf{Implementation:} Direct implementation via the convolution
  summation is \textbf{practically impossible} because it would require
  infinite memory, multiplications, and additions.
\item
  \textbf{Realization:} Instead of the convolution sum, IIR systems are
  more practically and efficiently implemented using \textbf{difference
  equations} and \textbf{recursive structures}. These structures use
  feedback loops to express the current output in terms of past output
  values as well as current and past inputs.
\end{itemize}

\subsubsection{Recursive and Nonrecursive Discrete-Time
Systems}\label{recursive-and-nonrecursive-discrete-time-systems}

\begin{itemize}
\item
  \textbf{Nonrecursive Systems}

  \begin{itemize}
  \tightlist
  \item
    \textbf{Definition:} A system is nonrecursive if its output \(y(n)\)
    at any time \(n\) depends only on the present and past values of the
    input, such that \(y(n) = F[x(n), x(n-1), \dots, x(n-M)]\).
  \item
    \textbf{Characteristics:} Causal \textbf{FIR (Finite-Duration
    Impulse Response) systems} are inherently nonrecursive.
  \item
    \textbf{Computation:} The output can be computed in any order (e.g.,
    \(y(200)\) can be found without first finding \(y(199)\)).
  \end{itemize}
\item
  \textbf{Recursive Systems}

  \begin{itemize}
  \tightlist
  \item
    \textbf{Definition:} A system is recursive if the current output
    \(y(n)\) is a function of past output values (e.g.,
    \(y(n-1), y(n-2)\)) in addition to current and past inputs.
  \item
    \textbf{Feedback Loop:} These systems are characterized by a
    \textbf{feedback loop} that contains at least one delay element,
    which is essential for the system to be practically realizable.
  \item
    \textbf{Computation:} Outputs must be computed sequentially in order
    (e.g., \(y(0), y(1), y(2), \dots\)).
  \item
    \textbf{Initial Conditions:} The response of a recursive system is
    not uniquely determined by the input alone; it also depends on the
    \textbf{initial conditions} (the state of the system's memory).
  \end{itemize}
\end{itemize}

\subsubsection{Implementation of Discrete-Time
Systems}\label{implementation-of-discrete-time-systems}

\begin{itemize}
\tightlist
\item
  \textbf{Linear Constant-Coefficient Difference Equations:} These are
  the primary means for implementing linear time-invariant (LTI)
  systems.
\item
  \textbf{Direct Form I:} This structure uses separate sets of delay
  elements (memory) for the input signal and the output signal. It can
  be viewed as a cascade of a nonrecursive system followed by a
  recursive system.
\item
  \textbf{Direct Form II:} By interchanging the order of the recursive
  and nonrecursive parts, the two sets of delays can be merged into a
  single set.

  \begin{itemize}
  \tightlist
  \item
    \textbf{Canonic Form:} Because it requires the \textbf{minimum
    number of memory locations} (\(max{M, N}\)), it is often referred to
    as the canonic form.
  \end{itemize}
\item
  \textbf{Modular Realization:} Higher-order systems are typically
  implemented by cascading \textbf{second-order sections} as basic
  building blocks.
\end{itemize}

\subsubsection{Correlation of Discrete-Time
Signals}\label{correlation-of-discrete-time-signals}

\begin{itemize}
\tightlist
\item
  \textbf{Purpose:} Correlation measures the degree to which two signals
  are \textbf{similar} to extract information for applications like
  radar, sonar, and digital communications.
\item
  \textbf{Crosscorrelation}

  \begin{itemize}
  \tightlist
  \item
    \textbf{Formula:} For two real sequences \(x(n)\) and \(y(n)\), the
    crosscorrelation is
    \[r_{xy}(l) = \sum_{n=-\infty}^{\infty} x(n)y(n-l)\]
  \item
    \textbf{Non-Commutative:} Unlike convolution, crosscorrelation is
    generally \textbf{not commutative}; changing the order of indices
    yields \(r_{xy}(l) = r_{yx}(-l)\), which is a folded version of the
    original sequence.
  \item
    \textbf{Relation to Convolution:} It can be computed using a
    convolution program by folding one of the sequences first:
    \(r_{xy}(l) = x(l) * y(-l)\).
  \end{itemize}
\item
  \textbf{Autocorrelation}

  \begin{itemize}
  \tightlist
  \item
    \textbf{Definition:} This occurs in the special case where a signal
    is correlated with itself (\(y(n) = x(n)\)).
  \item
    \textbf{Formula:}
    \[r_{xx}(l) = \sum_{n=-\infty}^{\infty} x(n)x(n-l)\]
  \item
    \textbf{Properties:} The autocorrelation sequence is an \textbf{even
    function} (\(r_{xx}(l) = r_{xx}(-l)\)) and attains its
    \textbf{maximum value at zero lag} (\(l=0\)).
  \item
    \textbf{Application:} It is used to identify periodicities in
    physical signals that may be heavily corrupted by random noise.
  \end{itemize}
\end{itemize}

\begin{quote}
\textbf{\emph{Example 6.2}}

Compute the \textbf{autocorrelation} of the signal

\[x(n) = a^n u(n), \quad 0 < a < 1\]

\emph{Solution}

\begin{enumerate}
\def\labelenumi{\arabic{enumi}.}
\tightlist
\item
  \textbf{Identify Signal Type:} \(x(n)\) is an infinite-duration
  signal; therefore, its autocorrelation sequence \(r_{xx}(l)\) will
  also have \textbf{infinite duration}.
\item
  \textbf{Remark:} We must split the solution into two cases, \(l > 0\)
  and \(l < 0\), because the \textbf{starting point of the overlap}
  between the signal \(x(n)\) and its shifted version \(x(n-l)\) changes
  based on the direction of the shift. Since the signal
  \(x(n) = a^n u(n)\) is \textbf{causal} (it is zero for \(n < 0\)), the
  product in the autocorrelation formula is only non-zero where both
  sequences are non-zero:
\end{enumerate}

\begin{itemize}
\tightlist
\item
  \textbf{Case 1 (\(l \ge 0\)):} Shifting \(x(n)\) to the right by \(l\)
  units means it starts at \(n = l\). The overlap with the original
  \(x(n)\) (which starts at \(n = 0\)) begins at \textbf{\(n = l\)}.
  Thus, the summation limits are from \(l\) to \(\infty\).
\item
  \textbf{Case 2 (\(l < 0\)):} Shifting \(x(n)\) to the left by \(l\)
  units means it starts at \(n = l\) (a negative value). The overlap
  with the original \(x(n)\) still begins at the original starting
  point, \textbf{\(n = 0\)}. Thus, the summation limits are from 0 to
  \(\infty\).
\end{itemize}

\begin{enumerate}
\def\labelenumi{\arabic{enumi}.}
\setcounter{enumi}{2}
\tightlist
\item
  \textbf{Evaluate for \(l > 0\):}

  \begin{itemize}
  \tightlist
  \item
    Apply the autocorrelation formula:
    \[r_{xx}(l) = \sum_{n=l}^{\infty} x(n)x(n-l)\]
  \item
    Substitute the signal:
    \[r_{xx}(l) = \sum_{n=l}^{\infty} a^n a^{n-l} = a^{-l} \sum_{n=l}^{\infty} (a^2)^n\]
  \item
    Change variables (\(k = n-l\)) and apply the geometric series
    formula: \[r_{xx}(l) = \frac{1}{1-a^2} a^l, \quad l > 0\]
  \end{itemize}
\item
  \textbf{Evaluate for \(l < 0\):}

  \begin{itemize}
  \tightlist
  \item
    Apply the formula:
    \[r_{xx}(l) = \sum_{n=0}^{\infty} x(n)x(n-l) = a^{-l} \sum_{n=0}^{\infty} (a^2)^n\]
  \item
    Apply the geometric series formula:
    \[r_{xx}(l) = \frac{1}{1-a^2} a^{-l}, \quad l < 0\]
  \end{itemize}
\item
  \textbf{Generalize the Result:} Since \(a^{-l} = a^{|l|}\) when \(l\)
  is negative, the two results can be combined into a single expression
  for all \(l\):
  \[r_{xx}(l) = \frac{1}{1-a^2} a^{|l|}, \quad -\infty < l < \infty\]
\end{enumerate}
\end{quote}

\begin{center}\rule{0.5\linewidth}{0.5pt}\end{center}

\subsection{The Z-Transform and Its Application to the Analysis of LTI
Systems}\label{the-z-transform-and-its-application-to-the-analysis-of-lti-systems}

\subsubsection{The Direct Z-Transform}\label{the-direct-z-transform}

The idea of the \textbf{z-transform} is to transform the time-domain
signal \(x(n)\) into the z-domain, which does not depend on \(n\). It is
defined as the following power series:

\[X(z) \equiv \sum_{n=-\infty}^{\infty} x(n)z^{-n}\]

where \(z\) is a complex variable. The inverse procedure of obtaining
\(x(n)\) from \(X(z)\) is called the \textbf{inverse z-transform}.

\textbf{Region of Convergence (ROC):} The ROC is the set of all values
of \(z\) for which \(X(z)\) attains a finite value. Any time a
z-transform is cited, its ROC should also be indicated, as a
discrete-time signal is only \textbf{uniquely determined} by the
combination of its z-transform \(X(z)\) and its ROC.

\begin{quote}
Calculate the z-transform of these signals:

\begin{itemize}
\item
  \textbf{Unit Sample Signal} \(x(n) = \delta[n]\):

  \begin{itemize}
  \tightlist
  \item
    \textbf{Calculation:} Since \(\delta(n)\) is \(1\) only at \(n=0\)
    and \(0\) elsewhere:
    \[X(z) = \sum_{n=-\infty}^{\infty} \delta(n)z^{-n} = \delta(0)z^{-0} = 1 \cdot 1 = 1\]
  \item
    \textbf{Result:} \(X(z) = 1\)
  \item
    \textbf{ROC:} Entire \(z\)-plane.
  \end{itemize}
\item
  \textbf{Unit Step Signal} \(x(n) = u[n]\):

  \begin{itemize}
  \tightlist
  \item
    \textbf{Calculation:} Since \(u(n) = 1\) for \(n \ge 0\):
    \[X(z) = \sum_{n=0}^{\infty} (1)z^{-n} = \sum_{n=0}^{\infty} (z^{-1})^n\]
    Applying the Geometric Series property
    \(\sum_{n=0}^{\infty} A^n = \frac{1}{1-A}\) where \(A = z^{-1}\):
    \[X(z) = \frac{1}{1-z^{-1}}\]
  \item
    \textbf{ROC:} Converges if \(|z^{-1}| < 1 \implies |z| > 1\).
  \end{itemize}
\item
  \textbf{Exponential Signal} \(x(n) = \alpha^n u(n)\):

  \begin{itemize}
  \tightlist
  \item
    \textbf{Calculation:}
    \[X(z) = \sum_{n=0}^{\infty} \alpha^n z^{-n} = \sum_{n=0}^{\infty} (\alpha z^{-1})^n\]
    Applying the Geometric Series property \(\frac{1}{1-A}\) where
    \(A = \alpha z^{-1}\): \[X(z) = \frac{1}{1-\alpha z^{-1}}\]
  \item
    \textbf{ROC:} Converges if
    \(|\alpha z^{-1}| < 1 \implies |z| > |\alpha|\).
  \end{itemize}
\item
  \textbf{Anticausal Exponential Signal} \(x(n) = -\alpha^n u(-n-1)\):

  \begin{itemize}
  \tightlist
  \item
    \textbf{Calculation:} This signal is non-zero only for \(n \le -1\).
    Let \(l = -n\):
    \[X(z) = \sum_{n=-\infty}^{-1} (-\alpha^n)z^{-n} = -\sum_{l=1}^{\infty} (\alpha^{-1}z)^l\]
    Applying the Geometric Series property
    \(\sum_{l=1}^{\infty} A^l = \frac{A}{1-A}\) where
    \(A = \alpha^{-1}z\):
    \[X(z) = -\frac{\alpha^{-1}z}{1-\alpha^{-1}z} = \frac{1}{1-\alpha z^{-1}}\]
  \item
    \textbf{ROC:} Converges if
    \(|\alpha^{-1}z| < 1 \implies |z| < |\alpha|\).
  \end{itemize}
\end{itemize}
\end{quote}

\textbf{Key properties:}

\begin{itemize}
\tightlist
\item
  \textbf{Linearity:} The transform of a linear combination of signals
  is the same linear combination of their transforms.
\item
  \textbf{Time Shifting:} Delaying a signal by \(k\) samples corresponds
  to multiplying its z-transform by \(z^{-k}\).
\item
  \textbf{Convolution:} Convolution in the time domain is equivalent to
  \textbf{multiplication} in the z-domain: \(Y(z) = H(z)X(z)\). This
  property greatly simplifies the analysis of LTI systems.
\end{itemize}

\begin{quote}
\textbf{\emph{Table 1 Characteristic Families of Signals with Their
Corresponding ROCs}}

\textbf{Finite-Duration Signals}

For signals that exist only over a specific, finite interval of time,
the ROC is generally the \textbf{entire \(z\)-plane}, with specific
points excluded depending on where the signal's non-zero samples are
located.

\begin{itemize}
\tightlist
\item
  \textbf{Causal:} The signal is zero for \(n < 0\).

  \begin{itemize}
  \tightlist
  \item
    \textbf{ROC:} Entire \(z\)-plane except \textbf{\(z = 0\)}. The
    origin is excluded because terms with \(z^{-k}\) (where \(k > 0\))
    become unbounded at \(z = 0\).
  \end{itemize}
\item
  \textbf{Anticausal:} The signal is zero for \(n > 0\).

  \begin{itemize}
  \tightlist
  \item
    \textbf{ROC:} Entire \(z\)-plane except \textbf{\(z = \infty\)}.
    \(z = \infty\) is excluded because terms with \(z^k\) (where
    \(k > 0\)) become unbounded there.
  \end{itemize}
\item
  \textbf{Two-sided:} The signal has non-zero samples for both positive
  and negative values of \(n\).

  \begin{itemize}
  \tightlist
  \item
    \textbf{ROC:} Entire \(z\)-plane except \textbf{\(z = 0\) and
    \(z = \infty\)}.
  \end{itemize}
\end{itemize}

\textbf{Infinite-Duration Signals}

For signals that continue indefinitely in one or both directions, the
ROC takes the form of circular regions or rings.

\begin{itemize}
\tightlist
\item
  \textbf{Causal:} These signals are zero for \(n < 0\) and continue to
  \(n = \infty\).

  \begin{itemize}
  \tightlist
  \item
    \textbf{ROC:} The \textbf{exterior of a circle} of radius \(r_2\)
    (\(|z| > r_2\)).
  \end{itemize}
\item
  \textbf{Anticausal:} These signals are zero for \(n \ge 0\) and
  continue to \(n = -\infty\).

  \begin{itemize}
  \tightlist
  \item
    \textbf{ROC:} The \textbf{interior of a circle} of radius \(r_1\)
    (\(|z| < r_1\)).
  \end{itemize}
\item
  \textbf{Two-sided:} These signals have infinite duration in both the
  positive and negative directions.

  \begin{itemize}
  \tightlist
  \item
    \textbf{ROC:} An \textbf{annular region (ring)} in the \(z\)-plane
    defined by \(r_2 < |z| < r_1\). If \(r_2 > r_1\), the two
    components' convergence regions do not overlap, and the
    \(z\)-transform does not exist for that signal.
  \end{itemize}
\end{itemize}
\end{quote}

\begin{quote}
\textbf{\emph{Table 2 Properties of the z-Transform}}

\textbf{Fundamental Properties}

\begin{itemize}
\tightlist
\item
  \textbf{Linearity:} The \(z\)-transform of a linear combination of
  signals is the same linear combination of their individual transforms:
  \(a_1x_1(n) + a_2x_2(n) \leftrightarrow a_1X_1(z) + a_2X_2(z)\). The
  resulting ROC is at least the \textbf{intersection} of the individual
  ROCs.
\item
  \textbf{Time Shifting:} Delaying or advancing a signal by \(k\)
  samples corresponds to multiplying \(X(z)\) by \(z^{-k}\):
  \(x(n - k) \leftrightarrow z^{-k}X(z)\). This is a critical property
  for solving \textbf{difference equations}.
\item
  \textbf{Scaling in the \(z\)-Domain:} Multiplying a signal by an
  exponential sequence \(a^n\) results in scaling the \(z\) variable:
  \(a^nx(n) \leftrightarrow X(a^{-1}z)\). This results in
  \textbf{expanding or shrinking} the \(z\)-plane and its ROC by a
  factor of \(|a|\).
\item
  \textbf{Time Reversal:} Folding a signal \(x(n)\) to \(x(-n)\)
  corresponds to replacing \(z\) with \(z^{-1}\):
  \(x(-n) \leftrightarrow X(z^{-1})\). The ROC is inverted
  (\(1/r_1 < |z| < 1/r_2\)).
\end{itemize}

\textbf{Signal Manipulation Properties}

\begin{itemize}
\tightlist
\item
  \textbf{Conjugation:} For complex signals,
  \(x^*(n) \leftrightarrow X^*(z^*)\).
\item
  \textbf{Real and Imaginary Parts:}
  \(Re\{x(n)\} \leftrightarrow \frac{1}{2}[X(z) + X^*(z^*)]\) and
  \(Im\{x(n)\} \leftrightarrow \frac{1}{2j}[X(z) - X^*(z^*)]\).
\item
  \textbf{Differentiation in the \(z\)-Domain:} Multiplying a signal by
  \(n\) is equivalent to differentiating its transform and multiplying
  by \(-z\): \(nx(n) \leftrightarrow -z \frac{dX(z)}{dz}\). This is
  useful for finding transforms of \textbf{ramp sequences}.
\end{itemize}

\textbf{System and Energy Properties}

\begin{itemize}
\tightlist
\item
  \textbf{Convolution:} This is arguably the most powerful property:
  convolution in the time domain is equivalent to
  \textbf{multiplication} in the \(z\)-domain:
  \(x_1(n) * x_2(n) \leftrightarrow X_1(z)X_2(z)\). This converts
  complex time-domain operations into simple algebraic ones.
\item
  \textbf{Correlation:} The \(z\)-transform of the crosscorrelation
  \(r_{x_1 x_2}(l)\) is \(X_1(z)X_2(z^{-1})\).
\item
  \textbf{Multiplication of Two Sequences:} Multiplying two signals in
  time corresponds to a \textbf{complex contour integral} (periodic
  convolution) in the \(z\)-domain.
\item
  \textbf{Parseval's Relation:} Relates the total energy of signals in
  the time domain to their representation in the \(z\)-domain via a
  contour integral.
\end{itemize}

\textbf{Initial Value Theorem}

If a signal \(x(n)\) is \textbf{causal} (zero for \(n < 0\)), the first
sample value \(x(0)\) can be found directly from \(X(z)\) by taking the
limit as \(z\) approaches infinity:
\textbf{\(x(0) = \lim_{z \to \infty} X(z)\)}.
\end{quote}

\begin{quote}
\textbf{\emph{Table 3 Some Common z-Transform Pairs}}

\textbf{Elementary Basic Signals}

\begin{itemize}
\tightlist
\item
  \textbf{Unit Sample \(\delta(n)\):} The transform is simply \(1\) for
  the entire \(z\)-plane. This is because the signal is only non-zero at
  \(n=0\).
\item
  \textbf{Unit Step \(u(n)\):} Expressed as
  \textbf{\(\frac{1}{1-z^{-1}}\)} with a Region of Convergence (ROC) of
  \textbf{\(|z| > 1\)}. This represents the sum of a geometric series
  for a causal signal.
\end{itemize}

\textbf{Exponential and Ramp Signals}

A critical takeaway from this section is that different time-domain
signals can result in the same algebraic \(X(z)\) expression; they are
only distinguished by their \textbf{ROC}.

\begin{itemize}
\tightlist
\item
  \textbf{Causal Exponential \(a^n u(n)\):}
  \(X(z) = \frac{1}{1-az^{-1}}\) with ROC \textbf{\(|z| > |a|\)}.
\item
  \textbf{Anticausal Exponential \(-a^n u(-n-1)\):} Shares the
  \textbf{same algebraic form} (\(\frac{1}{1-az^{-1}}\)) but has an ROC
  of \textbf{\(|z| < |a|\)}.
\item
  \textbf{Ramp-like Signals (\(na^n u(n)\) and \(-na^n u(-n-1)\)):}
  These result in \textbf{\(\frac{az^{-1}}{(1-az^{-1})^2}\)}. These are
  derived using the \textbf{differentiation property} (multiplying a
  signal by \(n\) in time corresponds to differentiation in the
  \(z\)-domain).
\end{itemize}

\textbf{Sinusoidal and Weighted Sinusoidal Signals}

These pairs represent oscillatory behavior and are derived using
\textbf{Euler's identity} to split the sine and cosine functions into
complex exponentials.

\begin{itemize}
\tightlist
\item
  \textbf{Standard Sine and Cosine:}

  \begin{itemize}
  \tightlist
  \item
    \textbf{\(\cos(\omega_0 n)u(n)\)}
    \(\leftrightarrow \frac{1-z^{-1}\cos\omega_0}{1-2z^{-1}\cos\omega_0 + z^{-2}}\).
  \item
    \textbf{\(\sin(\omega_0 n)u(n)\)}
    \(\leftrightarrow \frac{z^{-1}\sin\omega_0}{1-2z^{-1}\cos\omega_0 + z^{-2}}\).
  \item
    Both have an ROC of \textbf{\(|z| > 1\)}.
  \end{itemize}
\item
  \textbf{Exponentially Weighted Sine and Cosine (\(a^n \cos\dots\) and
  \(a^n \sin\dots\)):}

  \begin{itemize}
  \tightlist
  \item
    These follow the same structure as the standard sinusoids but
    include the scaling factor \textbf{\(a\)}.
  \item
    Their ROC is \textbf{\(|z| > |a|\)}, illustrating the
    \textbf{Scaling property} where multiplying by \(a^n\) in time
    scales the \(z\) variable.
  \end{itemize}
\end{itemize}
\end{quote}
